\chapter*{RESUMEN}
\addcontentsline{toc}{chapter}{RESUMEN}

La ampacidad de un cable eléctrico enterrado depende de las pérdidas internas y de la capacidad del entorno para evacuar calor. La humedad, los estratos, el lecho de asiento y el suelo de relleno producen una conductividad térmica espacialmente variable que puede modificar el campo de temperatura, la temperatura máxima del conductor y la corriente admisible. Esta investigación diseña y evalúa un artefacto computacional bidimensional basado en una red neuronal informada por física. El artefacto incorpora la ecuación estacionaria de conducción, las condiciones de contorno e interfaz y la conductividad espacial dentro del entrenamiento. La ciencia del diseño organiza la especificación, construcción, demostración, evaluación y comunicación del artefacto. La verificación considera soluciones analíticas o manufacturadas, referencias FEM convergentes y casos publicados; además, registra exactitud térmica, consistencia física, robustez y reproducibilidad.

\pendiente{Completar el resumen con los resultados cuantitativos principales, la conclusión central y el impacto estimado. Ajustar la extensión y las palabras clave después de cerrar el Capítulo III.}

\noindent\textbf{Palabras clave:} cable subterráneo; ampacidad; heterogeneidad térmica; conductividad del suelo; red neuronal informada por física; ciencia del diseño.

\chapter*{ABSTRACT}
\addcontentsline{toc}{chapter}{ABSTRACT}

The ampacity of an underground power cable depends on internal losses and on the ability of the surrounding environment to dissipate heat. Moisture, soil layers, bedding, and backfill create a spatially varying thermal conductivity that may alter the temperature field, the maximum conductor temperature, and the allowable current. This research designs and evaluates a two-dimensional computational artifact based on a physics-informed neural network. The artifact embeds the steady-state heat equation, boundary and interface conditions, and spatial conductivity into training. Design science organizes the specification, construction, demonstration, evaluation, and communication of the artifact. Verification uses analytical or manufactured solutions, converged FEM references, and published cases, while recording thermal accuracy, physical consistency, robustness, and reproducibility.

\pendiente{Traducir la versión final del resumen después de incorporar resultados, conclusión e impacto. Verificar que ambas versiones sean equivalentes.}

\noindent\textbf{Keywords:} underground cable; ampacity; thermal heterogeneity; soil thermal conductivity; physics-informed neural network; design science research.

